\chapter{进程同步}

\section{同步与互斥的概念}
\concept{并发进程的制约关系}
\begin{points}
  \item 多道程序系统中，由于资源共享或进程合作，并发进程之间形成相互制约关系。
  \item 间接相互制约：多个进程竞争同一资源，例如多个进程争用打印机。
  \item 直接相互制约：进程之间存在合作顺序，例如生产者生产后消费者才能消费。
  \item 进程同步的任务是协调这些关系，使并发进程能正确共享资源和相互合作，保证结果可再现。
\end{points}
\plain{同步这一章关心的不是“有没有并发”，而是“并发以后怎么别把共享数据搞坏”。}
\exam{常考直接/间接制约、为什么并发会导致执行结果不可再现。}

\concept{同步与互斥}
\begin{points}
  \item \term{同步}针对合作关系：多个进程在某些关键点上需要互相等待或按次序推进。
  \item \term{互斥}针对竞争关系：某个进程使用临界资源时，其他进程必须等待它释放资源。
  \item 同步强调“先后次序正确”；互斥强调“同一时刻只能一个进入”。
  \item 两者常同时出现：例如生产者-消费者问题中，访问缓冲区要互斥，放入/取出产品又要满足同步关系。
\end{points}
\plain{区分这两个词最稳的方法是看问题本质：如果是在管“谁先谁后”，偏同步；如果是在管“能不能同时进”，偏互斥。}
\exam{简答题经常直接问同步和互斥的区别，也可能给场景让你判断哪些关系属于同步、哪些属于互斥。}


\concept{与时间有关的错误/race condition}
\begin{points}
  \item 多个进程并发读写共享变量，执行次序不同可能导致不同结果，这种错误称为与时间有关的错误或竞争条件。
  \item 例子：订票系统中 A、B 同时读取剩余票数 $x=1$，都判断有票并出票，最终可能超卖。
  \item 根本原因：对共享变量的复合操作没有互斥保护。
  \item 解决思路：把读-改-写共享变量的代码段放入临界区，一次只允许一个进程进入。
\end{points}
\plain{竞争条件的关键不是“同时访问”这四个字，而是多个进程把一次完整更新拆成了几步，步骤交叉后就会把共享结果写乱。}
\exam{常考给一段并发读改写代码，让你解释为什么可能超卖、少减或结果不唯一。}


\concept{与时间有关的另一类错误：永远等待}
\begin{points}
  \item 与时间有关的错误不只会造成“结果不唯一”，还可能造成“该等的人没被正确唤醒”，从而永远等待。
  \item 典型情形是：进程刚判断完条件、还未来得及把自己放入等待队列，另一个进程就先释放资源并完成唤醒操作。
  \item 结果是释放动作已经发生，但等待进程还没真正登记到等待队列中，后续再也没人唤醒它。
  \item 这说明同步操作必须把“检查条件、入队等待、释放资源、唤醒他人”这些关键动作作为受控整体来设计。
\end{points}
\plain{有的并发错误不是把结果算错，而是把唤醒时机错过去，最后进程一直睡着没人再叫。}
\exam{常考“为什么简单 if/while + 等待队列操作也会出错”，答案要点是检查条件与入队/唤醒不是原子完成。}


\concept{临界资源与临界区}
\begin{points}
  \item 临界资源是一段时间内只允许一个进程使用的资源。
  \item 临界区是进程中访问临界资源的代码段。
  \item 进入区负责申请进入临界区；退出区负责释放临界区；剩余区是其他代码。
  \item 互斥要求：任意时刻最多一个进程在相关临界区内执行。
\end{points}
\plain{临界资源是“对象”，临界区是“访问这个对象的那段代码”；题目里真正要保护的是代码段，不是抽象名词本身。}
\exam{选择题很爱把临界资源和临界区混着考，答题时最好直接举“打印机/共享变量”和“访问它的代码段”来区分。}


\concept{临界区访问原则}
\begin{points}
  \item \key{空闲让进}：临界区空闲时，应允许一个请求进入的进程立即进入。
  \item \key{忙则等待}：临界区已有进程时，其他请求进程必须等待。
  \item \key{有限等待}：等待进入临界区的进程不能无限期等待，避免饥饿。
  \item \key{让权等待}：不能进入临界区的进程应释放处理机，不应长时间忙等浪费 CPU。
\end{points}
\plain{四条规则里最容易被忽略的是“让权等待”：进不去就该睡，不应一直空转耗 CPU，这也是后面记录型信号量优于忙等法的原因。}
\exam{常考把某个互斥方案和四条准则对照，判断它是否满足有限等待或让权等待。}


\concept{课本中的三原则表述}
\begin{points}
  \item 教材中常把临界区要求概括为：\key{互斥原则}、\key{推进原则}、\key{有限等待原则}。
  \item 互斥原则对应“同一时刻最多一个进程在临界区”。
  \item 推进原则强调当临界区空闲时，不能无限拖延某个请求进入的进程。
  \item 有限等待原则强调不能让某个进程无限期等下去。
\end{points}
\plain{不同资料有“四条规则”和“三条原则”两种说法，考试里看到哪种表述都要能对上同一组含义。}
\exam{常考把“空闲让进、忙则等待、有限等待、让权等待”和“互斥/推进/有限等待”对应起来。}

\section{互斥的实现方法}
\concept{软件方法}
\begin{points}
  \item 单标志法、双标志先检查法、双标志后检查法、Peterson 算法都是早期软件互斥方法。
  \item 软件方法的关键是用共享变量表达“谁想进”和“轮到谁进”。
  \item Peterson 算法能满足互斥、空闲让进和有限等待，但主要适合两个进程的理论模型。
  \item 软件方法缺点是复杂、容易出错、依赖忙等，难以推广。
\end{points}
\plain{软件法这节重点不在会背所有算法，而在看出它们共同缺陷：逻辑绕、推广差，而且很多方案本质上还在忙等。}
\exam{常考 Peterson 为什么理论上正确、却不适合现代通用系统实现；也常问软件法共同短板。}

\concept{Peterson 算法}
\begin{points}
  \item Peterson 算法面向两个进程，结合了\term{意愿标志} \code{flag[i]} 与\term{轮转变量} \code{turn}。
  \item 进程进入前先声明“我想进”（置 \code{flag[i]=true}），再把优先权让给对方（设置 \code{turn}）。
  \item 若对方也想进入，且当前 \code{turn} 指向对方，则本进程在 while 条件中等待；否则进入临界区。
  \item 退出临界区时把自己的 \code{flag[i]} 复位为 \code{false}。
\end{points}
\plain{Peterson 的直觉是“先表明意愿，再主动谦让一次”；如果对方此时也想进，就按 \code{turn} 决定谁先走，因此既避免同时进入，也避免两边永久互让。}
\exam{代码题常考 \code{flag} 和 \code{turn} 各表示什么，或问它为什么能满足互斥与有限等待。}

\concept{Dekker 算法}
\begin{points}
  \item Dekker 算法也是面向两个进程的软件互斥算法，结合了意愿标志 \code{flag[]} 和轮转变量 \code{turn}。
  \item 它的核心思想是：若双方都想进入临界区，则按 \code{turn} 决定谁暂时后退，让另一方先进入。
  \item 因而 Dekker 算法修正了“双方都置位后互相谦让到谁也进不去”的问题，是一个正确的软件互斥方案。
  \item 但其逻辑较复杂，通常更多作为理论过渡，后续常用更简洁的 Peterson 算法来讲解。
\end{points}
\plain{Dekker 可以理解成 Peterson 之前“更绕但能用”的版本：它已经把“想进”和“轮到谁”两类信息结合起来了，只是代码结构更复杂。}
\exam{若题目问 Dekker 和 Peterson 的共同思想，抓住“意愿标志 + 轮转变量”这条主线即可。}


\concept{硬件方法}
\begin{points}
  \item 关中断：单处理器中进入临界区前关中断，退出后开中断，可防止进程被中断切换。
  \item Test-and-Set 指令：原子地测试并设置锁变量。
  \item Swap 指令：原子地交换两个变量的值。
  \item 硬件方法可保证互斥，但若采用自旋锁，会造成忙等，不满足让权等待。
\end{points}
\plain{硬件法的价值是“提供原子性”，短板是很多实现仍让进程原地自旋，所以它解决了互斥，不一定解决了让权等待。}
\exam{常考为什么关中断只适合单处理器、为什么 Test-and-Set 能保互斥却仍可能浪费 CPU。}

\concept{关中断方法}
\begin{points}
  \item 在单处理器系统中，进程切换通常依赖时钟中断或 I/O 中断；进入临界区前先关中断，可避免当前进程被切走。
  \item 因而可用“关中断 $\rightarrow$ 执行临界区 $\rightarrow$ 开中断”的方式实现互斥。
  \item 它的主要缺点是会推迟中断响应、降低系统并发性，而且对多处理器无效。
  \item 由于若用户随意关中断可能导致系统失控，所以这种方法通常只适合内核中的短临界区。
\end{points}
\plain{关中断法的本质不是“上锁”，而是直接暂时禁止调度发生；所以它虽然简单，但代价是把系统响应能力一起压住了。}
\exam{比较题常问为什么关中断不适合多处理器，以及为什么不能把关中断权限交给普通用户进程。}

\concept{Test-and-Set 与 Swap 原子指令}
\begin{points}
  \item \term{Test-and-Set} 指令能原子地读出锁旧值并把锁置为占用，常写成 \code{while TS(\&lock);} 后进入临界区。
  \item \term{Swap} 指令能原子地交换共享锁变量与本地变量的值，也可据此反复尝试获取锁。
  \item 这类原子指令的共同点是：把“检查锁是否空闲”和“把锁改成占用”合并成不可分割的硬件操作。
  \item 它们适用于多处理器环境，但若失败后一直循环重试，就会形成忙等，不满足让权等待，且可能导致饥饿。
\end{points}
\plain{原子指令法真正解决的是“检查”和“上锁”之间不能被别人插进来；它修复了软件法最难保证的原子性问题，但没自动解决忙等。}
\exam{常考 TS/Swap 为什么能保证互斥，以及它们为什么仍可能浪费 CPU、不能保证有限等待。}


\concept{互斥锁机制}
\begin{points}
  \item 互斥锁 mutex lock 是操作系统提供的互斥访问工具，通常借助前面的原子硬件指令实现。
  \item 锁变量常用两种状态表示资源：0 表示资源可用（开锁），1 表示资源已被占用（关锁）。
  \item 进入临界区前执行 \code{lock}：若锁空闲则把锁置为 1 并进入；若锁已被占用，则继续等待。
  \item 退出临界区后执行 \code{unlock}，把锁恢复为 0，使其他等待者有机会进入。
  \item 若等待时一直反复检测锁状态，这种实现就是自旋锁；它有忙等开销，但在临界区很短、多处理器环境下常较实用。
\end{points}
\plain{互斥锁可以看成把“测试并上锁”“离开时解锁”包装成统一接口；真正的考点通常是它为什么简单好用、又为什么可能忙等浪费 CPU。}
\exam{PPT 把 mutex 和 spinlock 连着讲，选择题常考 0/1 两种状态、\code{lock}/\code{unlock} 的作用，以及自旋锁更适合“锁持有时间短”的场景。}


\section{信号量与 P/V 操作}
\concept{信号量定义}
\begin{points}
  \item 信号量 semaphore 由整数值 $s$ 和等待队列 $q$ 组成。
  \item 除初始化外，信号量的值只能通过 P 操作和 V 操作改变。
  \item 信号量 $s>0$ 时，表示可用资源数；$s=0$ 表示无可用资源但无人等待；$s<0$ 时，$|s|$ 表示等待队列中的进程数。
  \item P/V 操作必须是原子操作，执行过程中不能被其他进程打断。
\end{points}
\plain{定义题真正要抓的是“整数值 + 等待队列 + 只能被 P/V 改变 + 原子操作”，别只把它背成一个计数器。}
\exam{常考记录型信号量的组成，以及为什么 P/V 必须原子执行。}


\concept{P 操作}
\begin{points}
  \item P(S)：先执行 $s=s-1$。
  \item 若减完后 $s<0$，说明资源不够，当前进程阻塞，进入该信号量等待队列，并转调度程序。
  \item 若减完后 $s\ge 0$，说明成功获得资源，进程继续执行。
  \item 直觉：P 是“申请/占用一个资源”，先登记需求，再判断是否需要等待。
\end{points}
\plain{P 操作最容易写错的地方是“先减再判断”，不是先看够不够再减；这个顺序正是统一资源计数和等待队列语义的关键。}
\exam{常考手写 P 原语伪代码，或问为什么减完后小于 0 就表示当前进程必须阻塞。}


\concept{V 操作}
\begin{points}
  \item V(S)：先执行 $s=s+1$。
  \item 若加完后 $s\le 0$，说明仍有进程在等待队列中，唤醒一个等待进程并插入就绪队列。
  \item 若加完后 $s>0$，说明没有进程等待，只是增加了一个可用资源。
  \item 直觉：V 是“释放一个资源”。即使加完后 $s\le 0$，也要唤醒，因为负数表示有等待者；释放的资源应交给其中一个等待进程。
\end{points}
\plain{V 操作不是简单“把数加回去”，而是把释放出来的资源真正交还给等待队列里的某个进程，所以加完后若仍不大于 0 就得唤醒别人。}
\exam{常考为什么 V 后条件写成 `s <= 0` 而不是 `s < 0`，答题要结合“负数表示仍有人排队”等待来解释。}


\concept{互斥信号量取值范围}
\begin{points}
  \item 若 $N$ 个进程共享一个临界资源，互斥信号量初值通常设为 1。
  \item 当没有进程占用该资源时，信号量值可为 1。
  \item 当已有 1 个进程在临界区、其余部分进程等待时，信号量值会落到 0、-1、-2、...。
  \item 因此对一个临界资源的互斥信号量，其取值范围可理解为 $1,0,-1,\dots,-(N-1)$。
\end{points}
\plain{互斥信号量只有一份资源，所以初值是 1；后面一旦出现负数，本质上就是在数有多少人排队。}
\exam{常考“一个临界资源被 N 个进程共享时，互斥信号量初值和取值范围是什么”。}

\concept{整型信号量与记录型信号量}
\begin{points}
  \item 整型信号量只用整数表示资源数，P 操作中可能忙等。
  \item 记录型信号量增加等待队列，资源不足时让进程阻塞，满足让权等待。
  \item 考试中常默认 P/V 是记录型信号量，阻塞和唤醒由 OS 完成。
\end{points}
\plain{整型信号量和记录型信号量的根本差别不在“一个是 int 一个是 struct”，而在资源不足时是忙等还是阻塞让权。}
\exam{常考为什么记录型信号量比整型信号量更符合临界区的让权等待原则。}


\section{经典同步问题}
\concept{生产者-消费者问题}
\begin{points}
  \item 问题：多个生产者向有界缓冲区放产品，多个消费者从缓冲区取产品。
  \item 需要三个信号量：\code{mutex=1} 保护缓冲区互斥访问，\code{empty=n} 表示空缓冲区数，\code{full=0} 表示满缓冲区数。
  \item 生产者顺序：生产产品 $\rightarrow$ P(empty) $\rightarrow$ P(mutex) $\rightarrow$ 放入缓冲区 $\rightarrow$ V(mutex) $\rightarrow$ V(full)。
  \item 消费者顺序：P(full) $\rightarrow$ P(mutex) $\rightarrow$ 取出产品 $\rightarrow$ V(mutex) $\rightarrow$ V(empty) $\rightarrow$ 消费产品。
  \item 注意：先申请同步资源 empty/full，再申请互斥 mutex，避免拿着锁等待条件导致其他进程无法改变条件。
\end{points}
\plain{生产者消费者题最关键的顺序是先 P(empty/full) 再 P(mutex)，否则容易出现“拿着锁等条件”，把能改变条件的人也堵死。}
\exam{常考把 mutex 和 empty/full 的次序打乱后会不会死锁，或让你改正错误伪代码。}


\concept{P/V 题的规范写法}
\begin{points}
  \item 先定义每个信号量、初值以及它表示的含义。
  \item 伪代码中把每类进程的关键步骤写完整，不要只写零散几句 P/V。
  \item 写完后检查是否“有 P 就有对应的 V”，以及是否有人拿着互斥锁等待条件。
\end{points}
\plain{老师录音里反复强调：不会全写对也要规范写，定义清楚、关系写清楚，就能拿过程分。}
\exam{大题常按步骤给分，哪怕只分析出部分互斥/同步关系，也比空着强。}


\concept{读者-写者问题}
\begin{points}
  \item 多个读者可同时读共享文件，但写者必须互斥写；写者与读者也互斥。
  \item 读者优先方案中，第一个读者负责阻止写者，最后一个读者负责释放写者。
  \item 需要计数变量 \code{readcount} 记录当前读者数，并用互斥信号量保护它。
  \item 易错点：读者之间不是互斥读文件，而是互斥修改 \code{readcount}。
\end{points}
\plain{读者写者题最容易误判的是“读者共享读文件，但不共享改 `readcount` 这一步”；真正要互斥保护的是读者计数变量。}
\exam{常考第一个/最后一个读者分别负责什么，或者为什么读者之间不是完全没有互斥。}


\concept{前趋关系}
\begin{points}
  \item 若语句 $S_1$ 必须在 $S_2$ 之前完成，可设置初值为 0 的同步信号量 $S$。
  \item 前驱进程在完成 $S_1$ 后执行 V(S)。
  \item 后继进程在执行 $S_2$ 前先执行 P(S)。
  \item 前驱图较复杂时，可按每条前趋边设置同步信号量。
\end{points}
\plain{谁必须先做完，就让前面的人 V 一下，后面的人先 P 住等信号。}
\exam{常考把前驱图翻译成 P/V 代码。}


\concept{复杂前趋图的信号量设置方法}
\begin{points}
  \item 对复杂前趋图，常可按“每一条前趋边对应一个初值为 0 的同步信号量”来建模。
  \item 某进程若有多个直接前驱，则通常需要在开始前对多个同步信号量分别执行 P 操作。
  \item 某进程若有多个直接后继，则通常在结束后对多个同步信号量分别执行 V 操作。
  \item 这种方法虽然信号量较多，但结构清楚，考试中最稳妥。
\end{points}
\plain{多前驱就多次 P，多后继就多次 V，把图上的箭头一条条翻译掉，最不容易漏关系。}
\exam{常考六进程前趋图这类题，评分点通常就在“信号量初值设 0、前驱后继逐条翻译”上。}


\concept{哲学家进餐问题}
\begin{points}
  \item 五个哲学家围桌，每人左右各一只筷子，必须同时拿到两只筷子才能进餐。
  \item 若所有哲学家都先拿左筷子再等右筷子，可能形成死锁。
  \item 解决方法包括：最多允许四人同时取筷子；奇偶哲学家取筷子顺序相反；一次性申请两只筷子；使用管程。
  \item 核心直觉：避免每个进程“占有一个资源，同时等待另一个资源”。
\end{points}
\plain{哲学家进餐题不是单纯考代码套路，而是在追问“怎样打破占有一部分资源再等待另一部分资源”的死锁结构。}
\exam{常考判断某种取筷子策略是否会死锁，并说明它破坏了哪个死锁必要条件。}


\section{信号量集机制}
\concept{AND 型信号量}
\begin{points}
  \item AND 型信号量适用于进程同时需要多类资源、每类一个的情况。
  \item 基本思想：进程需要的多类资源一次性全部分配；只要有一个资源不满足，其他资源也不分配。
  \item 这样可以避免进程占有部分资源再等待其他资源，从而降低死锁风险。
  \item P 原语称为 SP/Swait，V 原语称为 SV/Ssignal。
\end{points}
\plain{它相当于“要么一起批，要么一个也别先拿”，避免拿到一半卡死。}
\exam{常考它为什么能缓解哲学家进餐这类“先占一部分再等另一部分”的死锁风险。}


\concept{一般信号量集}
\begin{points}
  \item 一般信号量集是 AND 型的扩展：一次可申请多类资源，每类资源可申请多个。
  \item 当某类资源低于下限或不能满足申请量时，本次原语不分配任何资源。
  \item 它把“多个 P 操作”合并成一个原子判断，避免中间状态。
\end{points}
\plain{关键不是新符号多，而是“多个资源请求合成一次原子判断”。}
\exam{常考它和基本记录型信号量、AND 型信号量的关系。}


\concept{信号量集的特殊退化形式}
\begin{points}
  \item 当信号量集写成 \code{SP(S,d,d)} 时，可理解为一次申请同类资源的 $d$ 个实例。
  \item 当写成 \code{SP(S,1,1)} 时，基本退化为普通记录型信号量。
  \item 当写成 \code{SP(S,1,0)} 时，可把它看成一种“开关型”控制：$S\ge 1$ 时允许进入，$S=0$ 时禁止进入。
\end{points}
\plain{这些特殊形式的价值在于帮你看出：信号量集并不是全新体系，而是把普通 P/V 和批量申请统一到一套写法里。}
\exam{常考信号量集和普通信号量的退化关系，尤其 \code{SP(S,1,1)} 为什么等价于基本信号量。}


\section{管程}
\concept{管程是什么}
\begin{points}
  \item 管程是一种高级同步机制，把共享数据、对共享数据的操作、同步条件封装在一个模块中。
  \item 管程保证任意时刻只有一个进程/线程在管程内部执行，因此天然提供互斥。
  \item 条件变量用于表达“某个条件不满足时等待，条件满足时被唤醒”。
  \item 典型操作包括 \code{wait} 和 \code{signal}。
\end{points}
\plain{这里先抓管程的“整体观”：共享数据、对数据的操作、同步条件被打包在一起，所以互斥不再靠程序员手动到处插 P/V。}
\exam{常考管程定义题，答题时最好同时写出“封装共享数据”和“任一时刻只允许一个进程进入”两层。}


\concept{管程的构成}
\begin{points}
  \item 管程通常由四部分组成：管程名、局部共享数据结构、对共享数据进行操作的一组过程、初始化共享数据的语句。
  \item 共享数据和相关操作被封装在同一抽象模块内，外部进程不能随意直接改动内部数据。
  \item 管程中的过程常可分为外部可调用过程和仅供管程内部使用的内部过程。
\end{points}
\plain{这节要抓“数据和操作被打包到一起”，它不像信号量那样只是给你一个原语，而更像带规则的抽象数据类型。}
\exam{PPT 单独列了管程构成，简答题里最好把“四部分”按顺序写出来，而不是只笼统说“有共享数据和过程”。}


\concept{管程的基本特性}
\begin{points}
  \item \term{安全性}：管程内的局部数据只能被管程内的过程访问。
  \item \term{共享性}：多个进程可通过调用同一个管程过程来共享受保护资源。
  \item \term{互斥性}：任意时刻至多一个进程真正进入管程内部执行。
  \item \term{透明性}：互斥进入等机制由语言/系统自动保证，程序员不必像写 P/V 那样手工到处安插同步操作。
\end{points}
\plain{这四条里最有管程味道的是“透明性”：程序员看到的是过程调用，互斥细节被语言和运行时隐藏起来。}
\exam{常考“为什么说管程比信号量更不容易写错”，答案通常要落到互斥自动化和透明性。}


\concept{入口等待队列与条件队列}
\begin{points}
  \item 由于管程一次只允许一个进程进入，所以当管程已被占用时，后来想进入的进程要先在\term{入口等待队列}排队。
  \item 若进程已经进入管程，但因为某个条件不满足而不能继续执行，则它会在对应的\term{条件变量等待队列}上睡眠。
  \item 入口等待队列反映的是“还没进门就在等互斥入口”；条件队列反映的是“已经进门，但因资源条件不满足而等待”。
  \item \code{wait} 会释放管程互斥权并转入相应条件队列；之后别的进程才能从入口队列进入管程继续执行。
\end{points}
\plain{这两个队列最容易混：一个是在门外排队等“进去”，一个是在屋里发现条件不够、主动让出管程后按原因排队等“继续”。}
\exam{PPT 把入口等待队列和条件变量分开讲了，简答题里若问管程内部等待机制，最好把这两层队列都写出来。}


\concept{管程与 P/V 的区别}
\begin{points}
  \item P/V 是低级同步原语，程序员需要自己安排每个 P、V 的位置，容易写错。
  \item 管程把共享变量和互斥规则封装起来，进入管程自动互斥，程序员主要表达条件等待。
  \item P/V 更像“手动锁门开门”；管程更像“把房间设计成一次只能进一个人，并提供排队等待条件”。
  \item 管程更结构化，适合高级语言；信号量更基础，适合描述底层同步。
\end{points}
\plain{这一节真正要比较的是同步风格：P/V 靠程序员手工排布原语，管程则把互斥和等待条件封进模块接口里，结构更清晰。}
\exam{常考比较题：为什么说管程比信号量更高级、更不容易出错，但底层实现仍可依赖 P/V。}


\concept{条件变量的 \code{wait}/\code{signal} 与 P/V 的区别}
\begin{points}
  \item 条件变量用于区分不同等待原因，每个条件变量对应一个条件队列。
  \item \code{x.wait} 会使当前进程在条件变量 $x$ 上等待，并释放管程互斥权。
  \item \code{x.signal} 唤醒 $x$ 队列中的一个等待进程；若队列为空，则相当于空操作。
  \item 与信号量不同，条件变量不做资源计数，没有累计值；它只维护等待队列。
\end{points}
\plain{条件变量更像“按原因排队”，不是“拿一个会累计加减的资源计数器”。}
\exam{常考为什么条件变量的 \code{signal} 在没人等待时不会“记账”，以及它和 V 操作的本质区别。}
